\chapter{GEN Ledger --- Generic Symbolic Primitives}
\label{chap:GEN_Ledger}

\section{Purpose}

GEN tags identify fundamental symbolic primitives, glyph-roots, structural
motifs, and pre-linguistic conceptual atoms used throughout the CME.

\section{Scope}

Includes:
\begin{itemize}
    \item symbolic primitives,
    \item glyph schemes,
    \item root morphisms,
    \item universal motifs,
    \item pre-semantic patterns.
\end{itemize}

\section{Schema}

\[
T_{\text{GEN}} = \langle key,\, gloss,\, bindings,\, constraints \rangle
\]

Examples:
\begin{itemize}
    \item \texttt{gen:glyph}
    \item \texttt{gen:primitive}
    \item \texttt{gen:root}
    \item \texttt{gen:form}
    \item \texttt{gen:signal}
\end{itemize}

\section{Class Binding Notes}

\begin{itemize}
    \item S-class dominant for structural meaning.
    \item O-class for mythic/symbolic narrative forms.
\end{itemize}

\section{Constraints}

\begin{itemize}
    \item GEN tags should be minimal and universal.
    \item Avoid redundancy with ONTO unless structural.
    \item GEN tags must remain substrate-independent.
\end{itemize}
